\section{动物保健：周期中唯一稳定的“卖水人”}
\label{sec:vet}

动物保健板块 \num{9} 家公司（全部归入申万三级“动物保健Ⅲ”），
流通市值合计约 \num{481} 亿元，2025 年仅 \num{2} 家亏损，
\hl{ROE 中位数 \num{7.10}\%，为农业全部子行业最高}，
资产负债率超过 \num{65}\% 的只有 \num{1} 家。

\begingroup\scriptsize\setlength{\tabcolsep}{2pt}
\begin{longtable}{@{}llrrrrrrr@{}}
\caption{农业上市公司分行业明细（动物保健Ⅲ、动物保健；财务为 2026 年半年报，ROE 未年化，公告计数窗口 2023-09---2026-09）}\label{tab:comp-vet}\\
\toprule
代码 & 公司简称 & 市值（亿元） & 营收（亿元） & 同比（\%） & 净利（亿元） & ROE（\%） & 再融资公告 & 重组公告 \\
\midrule
\endfirsthead
\multicolumn{9}{l}{\small（续）}\\
\toprule
代码 & 公司简称 & 市值（亿元） & 营收（亿元） & 同比（\%） & 净利（亿元） & ROE（\%） & 再融资公告 & 重组公告 \\
\midrule
\endhead
\bottomrule
\endlastfoot
600201 & 生物股份 & 125.4 & 7.6 & 22.1 & 1.04 & 1.9 & 43 & 0 \\
600195 & 中牧股份 & 67.8 & 31.7 & 13.5 & 0.50 & 0.9 & 3 & 2 \\
688526 & 科前生物 & 56.8 & 3.7 & -23.3 & 1.16 & 2.8 & 0 & 0 \\
300119 & 瑞普生物 & 56.4 & 17.2 & 0.8 & 2.04 & 4.3 & 4 & 13 \\
300871 & 回盛生物 & 42.2 & 8.6 & 5.2 & 2.06 & 9.1 & 46 & 2 \\
002688 & 金河生物 & 38.7 & 14.5 & 4.1 & 1.31 & 5.7 & 26 & 1 \\
603566 & 普莱柯 & 35.3 & 4.9 & -12.2 & 0.77 & 3.0 & 4 & 2 \\
688098 & 申联生物 & 34.4 & 1.1 & -8.1 & -0.19 & -1.4 & 3 & 2 \\
002868 & 绿康生化 & 24.1 & 2.5 & -15.7 & 0.07 & 10.9 & 8 & 23 \\
\end{longtable}
\endgroup


\subsection{为什么动保能在下行周期中盈利}
\label{sec:vet-mechanism}

动保的需求由\hl{存栏量而非价格}决定：即使猪价跌破成本线，
养殖场也不会减少疫苗接种与保健投入，因为疫病造成的损失远大于药费。
这使动保企业的收入与生猪、家禽存栏量挂钩，
而成本端（原料药、发酵产能）与农产品价格无关。
\emph{在农业下行周期中，动保是唯一同时具备“需求稳定 + 成本独立”特征的环节}。

板块盈利结构印证了这一点：\textbf{科前生物}（\num{4.19} 亿元、ROE \num{10.3}\%）、
\textbf{瑞普生物}（\num{4.01} 亿元、ROE \num{8.8}\%）、
\textbf{回盛生物}（\num{2.53} 亿元、ROE \num{13.4}\%）三家合计利润 \num{10.7} 亿元，
占板块 \num{14.0} 亿元利润总额的 \num{76.7}\%。
\hl{科前生物与普莱柯以猪用疫苗为主，其客户正是牧原、温氏等逆势扩张的龙头}——
龙头企业在周期底部提升出栏量，反而增加了对疫苗的需求，
这是动保与养殖景气度“脱钩”的微观基础。

板块内唯一的实质性亏损是 \textbf{绿康生化}（\num{-1.39} 亿元），
其业务重心已转向微生物发酵与兽药原料药，\hl{不再是纯粹的动保公司}；
\textbf{申联生物}亏损 \num{0.20} 亿元，主要受口蹄疫疫苗招标价格下行影响。
即便是处于价格竞争最激烈的口蹄疫疫苗市场，
板块的亏损面（22.2\%）也远低于饲料（53.3\%）和养殖（40.9\%）。

\subsection{后非洲猪瘟时代的结构性变化}
\label{sec:vet-africa}

非洲猪瘟疫情（2018 年 8 月起，见第 \ref{sec:hog-shock} 节）以来，
养殖行业的生物安全标准被永久性抬高：
\hl{“低猪价 + 高防疫投入”成为新常态}。
这带来两个结构性变化：

\begin{enumerate}[leftmargin=1.5em, itemsep=2pt]
  \item \textbf{需求从“治疗”转向“预防”}：疫苗与诊断试剂占比提升，
        化药占比下降。板块内以疫苗为主的公司（科前生物、普莱柯、生物股份）
        ROE 均高于以化药为主的公司。
  \item \textbf{行业集中度提升}：中小养殖场因生物安全投入不足而退出，
        动保客户结构随之向大型集团场集中。
        \hl{这使动保企业获得了对下游的议价能力}，
        在 \num{2024}---\num{2025} 年养殖端普遍亏损时仍能维持毛利率。
\end{enumerate}

从融资行为看，动保板块近三年人均“再融资相关公告” \num{15.2} 条（中位数 \num{4} 条），
明显低于养殖（\num{19.2} / \num{11}）与饲料（\num{17.6} / \num{9}）。
\hl{盈利稳定 + 低负债 + 低融资依赖，使动保成为农业板块中资产负债表最健康的子行业}。
其中 \textbf{绿康生化}（\num{23} 条重组并购公告）与 \textbf{瑞普生物}（\num{13} 条）
的公告密度显著高于同业，反映出这两家在主动进行外延式扩张或业务转型。

\begin{keybox}[动物保健的画像结论]
\normalsize
\textbf{① 唯一的“逆周期”子行业}：ROE 中位数 \num{7.10}\%、
亏损面 \num{22.2}\%、高负债公司仅 \num{1} 家，三项指标均为农业最优。\par\vspace{3pt}
\textbf{② 需求的锚是存栏量，不是价格}：龙头企业在下行周期扩张产能，
反而推升了疫苗需求，这是动保与猪周期“脱钩”的机制。\par\vspace{3pt}
\textbf{③ 关注外延扩张的两家}：绿康生化转型原料药、瑞普生物持续并购，
\hl{它们代表动保行业从“单一疫苗”走向“疫苗 + 化药 + 诊断”平台化的路径}。
\end{keybox}
